% Generated from 00-abstract-en.md; edit Markdown, then rebuild.
Coding agents can read source files, edit code, and execute commands through tool use. However, producing a code change does not establish that a defect has been repaired. Repository-level repair also requires execution isolation, explicit acceptance criteria, failure feedback, and resource control. Multi-agent collaboration separates planning, implementation, and independent inspection, while introducing additional context and coordination costs.

This thesis designs and implements a multi-agent collaboration and verification system for code repair on top of an existing Python terminal coding agent. Coordinator, worker, and verifier roles operate with independent contexts and sequential handoffs. Changes are made in a disposable workspace, commands run in restricted Docker containers, and caller-specified tests execute in a separate verification copy. Execution evidence guides bounded retries under model-call, token, attempt, and time limits. The system exports candidate patches and execution reports. Skills provide reusable repair procedures, while project-scoped memory stores repair experiences with provenance, source fingerprints, and expiration policies.

The verification record reports 160 passing automated tests, including existing regression tests. Ten artificial defects were injected into a pinned version of python-dotenv: two development tasks and eight held-out evaluation tasks. Across 24 real model runs under the same budget settings, the single-role execution-feedback strategy, multi-role review-feedback strategy, and multi-role verification-feedback strategy completed acceptance in eight, five, and four cases, respectively. No false acceptances were observed in any group. Multi-role failures resulted from budget exhaustion; one candidate in the verification-feedback group passed additional tests without completing the workflow within budget.

Functional tests validate the implemented isolation, acceptance, retry, and evidence-delivery mechanisms. The repair experiments do not demonstrate a success-rate advantage for multi-role collaboration on this small task set. Engineering evaluation of coding agents should consider acceptance evidence, coordination costs, and stopping conditions together.

\textbf{Keywords:} coding agent; multi-agent collaboration; verification feedback; container isolation; experience memory
